\section{UVM Tests and Sequences}

The verification environment utilizes three different UVM Tests to verify the DUT. Each test is a wrapper that instantiates the UVM environment, raises an objection to prevent the simulation from ending, starts a specific sequence on the sequencer of the agent, and then drops the objection once the sequence finishes.

\subsection{Smoke Test (\texttt{gcd\_smoke\_test})}
The Smoke Test is designed to verify the fundamental connectivity of the testbench and ensure the DUT can successfully complete a basic transaction without hanging. 

The Smoke Test executes the \texttt{gcd\_smoke\_seq}, which overrides the default sequence item constraints to inject a single, hardcoded transaction. The operands are hardcoded to $A = 15$ and $B = 5$, and the \texttt{out\_ready\_delay} is explicitly forced to 0. This represents the simplest possible handshake scenario, to confirm that the interface protocols, reset behaviors, and FSM transitions function correctly before introducing randomness or complex timing.

\subsection{Directed Edge Case Test (\texttt{gcd\_directed\_test})}
To ensure the DUT properly handles boundary conditions and early exit scenarios (with only 1 cycle latency), the \texttt{gcd\_directed\_test} executes the \texttt{gcd\_directed\_edge\_case\_seq}. 

Instead of relying on random probability to hit specific mathematical extremes, this sequence explicitly utilizes a custom \texttt{send\_directed\_item()} task to add targeted operand combinations. The following edge cases are verified:
\begin{itemize}
    \item \textbf{Zero Operands:} $\gcd(0,0)$, $\gcd(A,0)$, and $\gcd(0,B)$. These test the immediate early-exit transition from \texttt{IDLE} directly to \texttt{DONE}.
    \item \textbf{Equal Operands:} $\gcd(A,A)$. This verifies the immediate equality check bypassing the iterative subtraction loop.
    \item \textbf{Maximum Boundaries:} Combinations utilizing $(2^{WIDTH}) - 1$. This stresses the datapath width limits and verifies the combinatorial subtraction logic under maximum bit-toggle conditions.
\end{itemize}

\begin{minted}[frame=single,framesep=10pt,breaklines]{SystemVerilog}
task body();
    `uvm_info("SEQ", "Executing directed edge case sequence", UVM_LOW)

    send_directed_item(0, 0);                 // GCD(0, 0)
    send_directed_item(18, 0);                // GCD(A != 0, B = 0)
    send_directed_item(0, 24);                // GCD(A = 0, B != 0)
    send_directed_item(27, 27);               // GCD(A = B)
    
    // Max values (WIDTH-1)
    send_directed_item((1 << WIDTH) - 1, 5);  // GCD(MAX, B)
    send_directed_item(5, (1 << WIDTH) - 1);  // GCD(A, MAX)
    send_directed_item((1 << WIDTH) - 1, (1 << WIDTH) - 1); // GCD(MAX, MAX)

    send_directed_item(0, (1 << WIDTH) - 1); // GCD(0, MAX)
    send_directed_item((1 << WIDTH) - 1, 0); // GCD(MAX, 0)
endtask
\end{minted}

\subsection{Constrained Random Test (\texttt{gcd\_random\_test})}
The main part of the functional verification is performed by the \texttt{gcd\_random\_test}, which runs the \texttt{gcd\_random\_seq}. This test stresses the general iterative subtractive algorithm ($1+N$ cycles latency) and the output stalling logic.

\begin{minted}[frame=single,framesep=10pt,breaklines]{SystemVerilog}
class gcd_random_seq extends uvm_sequence #(gcd_sequence_item);
    `uvm_object_utils(gcd_random_seq)

    rand int number_transactions;

    constraint c_number_of_transactions {
        number_transactions inside {[250:500]}; // rand between 250 and 500
    }

    function new(string name = "gcd_random_seq");
        super.new(name);
    endfunction

    task body();
        if (!randomize()) begin
            `uvm_error("SEQ", "Randomization failed for number of transactions")
        end
        
        `uvm_info("SEQ", $sformatf("Executing Random Sequence with %0d transactions", number_transactions), UVM_LOW)

        for (int i = 0; i < number_transactions; i++) begin
            req = gcd_sequence_item::type_id::create("req");
            $display("Randomizing transaction %0d/%0d", i+1, number_transactions);
            start_item(req);
            if (!req.randomize()) begin
                `uvm_error("SEQ", "Randomization failed for transaction")
            end
            $display("Finishing transactions");
            finish_item(req);
            $display("Transactions finished");
        end
    endtask
endclass
\end{minted}

The sequence randomizes the \texttt{number\_transactions} variable to execute between 250 and 500 transactions. During this sequence, the \texttt{out\_ready\_delay} parameter within the sequence items frequently forces the driver to stall the output handshake, in order to test the \texttt{DONE} state stability requirement and the concurrent assertions monitoring it.